\documentclass[11pt]{scrartcl} % Font size

\usepackage{amsmath, amsfonts, amsthm, amssymb} % Math packages

\usepackage{pdfpages} %add pdf pages

\usepackage{float} %figures in place

\usepackage{physics,mathtools} %extra math

\usepackage{listings} % Code listings, with syntax highlighting

\usepackage{xcolor, soul}

\usepackage[english]{babel} % English language hyphenation

%\usepackage[xetex]{graphicx}
\usepackage{graphicx} % Required for inserting images
\graphicspath{{Figures/}{./}} % Specifies where to look for included images (trailing slash required)

\usepackage{booktabs} % Required for better horizontal rules in tables

\numberwithin{equation}{section} % Number equations within sections (i.e. 1.1, 1.2, 2.1, 2.2 instead of 1, 2, 3, 4)
\numberwithin{figure}{section} % Number figures within sections (i.e. 1.1, 1.2, 2.1, 2.2 instead of 1, 2, 3, 4)
\numberwithin{table}{section} % Number tables within sections (i.e. 1.1, 1.2, 2.1, 2.2 instead of 1, 2, 3, 4)

\setlength\parindent{0pt} % Removes all indentation from paragraphs

\usepackage{enumitem} % Required for list customisation
\setlist{noitemsep} % No spacing between list items


%----------------------------------------------------------------------------------------
%   HYPER LINKS
%----------------------------------------------------------------------------------------
\usepackage{hyperref} %for hyperlinks
\hypersetup{
    colorlinks=true,
    linkcolor=blue,
    filecolor=magenta,
    urlcolor=cyan,
}
\usepackage{extarrows} %xleftrightarrow

%----------------------------------------------------------------------------------------
%	DOCUMENT MARGINS
%----------------------------------------------------------------------------------------

\usepackage{geometry} % Required for adjusting page dimensions and margins

\geometry{
	paper=a4paper, % Paper size, change to letterpaper for US letter size
	top=2.5cm, % Top margin
	bottom=3cm, % Bottom margin
	left=3cm, % Left margin
	right=3cm, % Right margin
	headheight=0.75cm, % Header height
	footskip=1.5cm, % Space from the bottom margin to the baseline of the footer
	headsep=0.75cm, % Space from the top margin to the baseline of the header
	%showframe, % Uncomment to show how the type block is set on the page
}

%----------------------------------------------------------------------------------------
%	FONTS
%----------------------------------------------------------------------------------------

\usepackage[utf8]{inputenc} % Required for inputting international characters
\usepackage[T1]{fontenc} % Use 8-bit encoding

% \usepackage{fourier} % Use the Adobe Utopia font for the document

%----------------------------------------------------------------------------------------
%	SECTION TITLES
%----------------------------------------------------------------------------------------

\usepackage{sectsty} % Allows customising section commands
\sectionfont{\vspace{6pt}\centering\normalfont\scshape} % \section{} styling
\subsectionfont{\normalfont\bfseries\centering} % \subsection{} styling
\subsubsectionfont{\normalfont\itshape\centering} % \subsubsection{} styling
\paragraphfont{\normalfont\scshape} % \paragraph{} styling

%----------------------------------------------------------------------------------------
%	HEADERS AND FOOTERS
%----------------------------------------------------------------------------------------

\usepackage{scrlayer-scrpage} % Required for customising headers and footers

\ohead*{} % Right header
\ihead*{} % Left header
\chead*{} % Centre header

\ofoot*{} % Right footer
\ifoot*{} % Left footer
\cfoot*{\pagemark} % Centre footer
 % Include the file specifying the document structure and custom commands
% LaTeX settings for MATLAB code listings
% based on Ted Pavlic's settings in http://links.tedpavlic.com/ascii/homework_new_tex.ascii
\usepackage{listings}
% \usepackage[usenames,dvipsnames]{color}
\usepackage{color}

% This is the color used for MATLAB comments below
\definecolor{MyDarkGreen}{rgb}{0.0,0.4,0.0}

% For faster processing, load Matlab syntax for listings
\lstloadlanguages{Matlab}%
\lstset{language=Matlab,                        % Use MATLAB
        frame=single,                           % Single frame around code
        basicstyle=\scriptsize\ttfamily,        % Use small true type font
        keywordstyle=[1]\color{blue}\bfseries,  % MATLAB functions bold and blue
        keywordstyle=[2]\color{purple},         % MATLAB function arguments purple
        keywordstyle=[3]\color{blue}\underbar,  % User functions underlined and blue
        identifierstyle=,                       % Nothing special about identifiers
                                                % Comments small dark green courier
        commentstyle=\usefont{T1}{pcr}{m}{sl}\color{MyDarkGreen}\small,
        stringstyle=\color{purple},             % Strings are purple
        showstringspaces=false,                 % Don't put marks in string spaces
        tabsize=3,                              % 5 spaces per tab
        %
        %%% Put standard MATLAB functions not included in the default
        %%% language here
        morekeywords={xlim,ylim,var,alpha,factorial,poissrnd,normpdf,normcdf,imresize},
        %
        %%% Put MATLAB function parameters here
        morekeywords=[2]{on, off, interp},
        %
        %%% Put user defined functions here
        morekeywords=[3]{FindESS, homework_example, my_imfilter},
        %
        morecomment=[l][\color{blue}]{...},     % Line continuation (...) like blue comment
        numbers=left,                           % Line numbers on left
        firstnumber=1,                          % Line numbers start with line 1
        numberstyle=\tiny\color{blue},          % Line numbers are blue
        stepnumber=1                            % Line numbers go in steps of 5
        }

% Includes a MATLAB script.
% The first parameter is the label, which also is the name of the script
%   without the .m.
% The second parameter is the optional caption.
\newcommand{\matlabscript}[2]
  {\begin{itemize}\item[]\lstinputlisting[caption=#2,label=#1]{#1.m}\end{itemize}}


\usepackage{fontspec}
\setmainfont{Tinos Nerd Font} %nice font for english and greek

\begin{document}

\begin{titlepage}
    \centering
    \includegraphics[width=0.9\textwidth]{logo_tuc.png}\par\vspace{1cm}
    \normalfont\normalsize
    \textsc{\textcolor[rgb]{0.66, 0.09, 0.19}{Ηλεκτρολόγων Μηχανικών και Μηχανικών Υπολογιστών}}\\ % Your university, school and/or department name(s)
    \vspace{25pt} % Whitespace
    %\textcolor[rgb]{0.66, 0.09, 0.19}{\rule{\linewidth}{0.5pt}}
    \rule{\linewidth}{0.5pt}\\ % Thin top horizontal rule
    \vspace{20pt} % Whitespace
    {\Huge Μηχανική Όραση}\\ % The assignment title

    {\huge Αναφορά Δεύτερου Project}\\ % The assignment title
    \vspace{12pt} % Whitespace
    \rule{\linewidth}{2pt}\\ % Thick bottom horizontal rule
    \vspace{12pt} % Whitespace
    \vspace{2cm}

    {\LARGE{Τσιαούσης Χρήστος \hfill Μητσάκης Χαράλαμπος}
        \par
        \texttt{2016030017 \hfill 2006030075}
        \par
    }

    \vfill
    Διδάσκουσα

    Ε. Δούτση

    \vfill

% Bottom of the page
    {\large \today\par}
\end{titlepage}

\newpage
% \tableofcontents


\section{Εισαγωγή}

\section{Υλοποίηση}

\subsection{Δημιουργία φίλτρων και scale-space}

Απο το paper του Lowe γνωρίζουμε ότι το scale-space μπορεί να υπολογιστεί με συνέλιξη της εικόνας με τη διαφορά δύο gaussian φίλτρων ως εξής:

\[D(x,y,\sigma) = (G(x,y,k\sigma) − G(x,y,\sigma)) * I(x,y)\]

Όμως η διαφορά δύο gaussian φίλτρων μπορεί να προσεγγυστεί απο το laplacian of gaussian ως εξής:

\[G(x,y,k\sigma) − G(x,y,\sigma) \approx (k−1)\sigma^2 \nabla^2 G\]

Άρα:

\[D(x,y,\sigma) = ((k−1)\sigma^2 \nabla^2 G) * I(x,y)\]

Επομένως δημιουργούμε φίλτρα με περιττές διαστάσεις ως εξής:

\begin{verbatim}
for i = 1:log_scales_per_octave
  k_power = i-1;
  sigma_prime = k^k_power * sigma;
  log_filters{i} = (k - 1) * sigma_prime^2 * ...
    fspecial('log', floor(4*sigma_prime) * 2 + 1, sigma_prime);
end
\end{verbatim}

Αντί να δημιουργούμε φίλτρα για όλα τα scales χωρίζουμε το scale-space σε οκτάβες.

Η πρώτη εικόνα κάθε οκτάβας υπολογίζεται με υποδειγματοληψία της προηγούμενης οκτάβας στο $1/2$ του μεγέθους, εκτός απο την πρώτη οκτάβα που δημιουργείται με διπλασιασμό της αρχικής εικόνας με linar interpolation.

Κάθε οκτάβα αντιστοιχεί σε διπλασιασμό του $\sigma$, αλλα είναι πιο αποδοτικό, αντί να συνεχίζουμε να αυξάνουμε το $\sigma$, να κάνουμε downsample την εικόνα και να επαναχρησιμοποιούμε τα φίλτρα.

Έστω $s$ ο αριθμός των scales ανα οκτάβα στα οποία κάνουμε αναζήτηση extrema.
Κάθε οκτάβα περιέχει $s+2$ εικόνες. Γιαυτό δημιουργούμε $s+2$ φίλτρα τα οποία επαναχρησιμοποιούμε σε κάθε οκτάβα.
Οι 2 επιπλέον εικόνες υπάρχουν για να μπορέσουμε να κάνουμε αναζήτηση των extrema σε $s$ εικόνες γιατί για την αναζήτηση χρειαζόμαστε το επόμενο και το προηγούμενο scale όπως εξηγούμε παρακάτω (Εύρεση extrema).

Το $k$ υπολογίζεται απο το $s$ ως εξής: $k = 2^{1/s}$ γιατί κάθε οκτάβα αντιστοιχεί σε διπλασιασμό του $\sigma$ και η οκτάβα χωρίζεται σε $s$ διαστήματα.

Επιλέξαμε τις τιμές $\sigma = 1.6$, $k = 2^{1/3}$, $s = 3$ που είναι οι προτεινόμενες τιμές απο τον Lowe.
Επίσης επιλέξαμε $3$ οκτάβες, άρα έχουμε $n = 3 \cdot 5 = 15$ επίπεδα στα οποία κάναμε αναζήτηση extrema.

Το scale-space αποθηκεύεται σε ένα cell array με διαστάσεις $num\_of\_octaves \times (s+2)$ και κάθε στοιχείο του cell array περιέχει μια εικόνα.
\begin{verbatim}
scale_space = cell(num_of_octaves, s+2);
\end{verbatim}

\subsection{Εύρεση extrema}

Για να βρούμε τα τοπικά maxima και minima του scale-space (υποψήφια keypoints) συγκρίνουμε κάθε σημείο του scale-space με τα 26 γειτονικά του (8 απο το ίδιο scale + 9 απο το κατώτερο + 9 απο το ανώτερο scale).

Απορρίπτουμε υποψήφια keypoints με τιμές κοντά στο 0 κάτω απο ένα threshold.

Επίσης απορρίπτουμε υποψήφια keypoints στα edges ελέγχοντας αν ισχύει η σχέση (για $r = 10$):
\[\frac{Tr(H)^2}{Det(H)} < \frac{(r+1)^2}{r}\]

\section{Αποτελέσματα}

\section{Συμπεράσματα}

\end{document}
